\section{Methods}
\subsection{Study design and reference categories}
This retrospective study evaluated WSI-level classification of recorded salivary gland MALT lymphoma and SS diagnoses. MALT was the positive class and SS the negative class. An SS medical history in a patient with MALT was a separate clinical attribute, not a reason to assign that slide to the SS class. Training used slide labels without cell-, patch- or lesion-level supervision. Clinical sources, dates, patient--specimen--slide linkage and case characteristics are reserved in Table~\ref{tab:cohort}. The source records identify Peking University School and Hospital of Stomatology for the internal material and Peking University Third Hospital for an external series; the complete provenance of the separate HE-directory material remains a field for the hospital records.

\textbf{Reference standard.} \pending{H02: hospital-held diagnostic criteria, reference-review procedure, reviewer roles, ancillary tests and adjudication.} The exact criteria used in the source records will be stated here, rather than inferred from general lymphoma or SS classification guidance \citep{ref3,ref4,ref17}.

\textbf{Eligibility and selection.} \pending{H03: clinical inclusion/exclusion rules, selection personnel and chronology, reasons for four prediction failures and 55 non-inclusions after successful prediction.} The recorded archive, successful inference and final analysis sets were retained as distinct stages. The analysis did not redefine cohort membership. Reporting was structured with reference to STARD-AI; unresolved items are listed in the accompanying completion checklist \citep{ref20}.

\subsection{Image preparation and data partition}
TRIDENT with GrandQC was used for tissue segmentation and patch extraction \citep{ref12}. The 193 internal and 168 analyzed external WSI headers recorded a Hamamatsu C9600-12 scanner, a $40\times$ source-lens setting and approximately $0.225708\,\mu\mathrm{m}$ per level-0 pixel. Non-overlapping $512\times512$-pixel patches were extracted at target $20\times$ magnification, corresponding to $1,024\times1,024$ level-0 pixels. The frozen UNI2-h encoder generated 1,536-dimensional features \citep{ref8,uni2card}. Its archived preprocessing resized the shorter image side to 224 pixels, applied a $224\times224$ center crop and ImageNet channel normalization. \pending{T01: exact encoder-weight checksum and downloaded revision.} Scanner-header details and the implementation record are provided in Supplementary Methods.

The 193-slide UNI2-h dataset was divided by label-stratified random sampling with seed 42 and a validation fraction of 0.25, yielding 144 training and 49 validation slides. The existing author-reviewed partition was reused in all controlled experiments. Each neural input contained 512 patch vectors, sampled without replacement when possible and with replacement otherwise. Training inputs varied across epochs or loading events; reference evaluation bags were deterministic. All reported counts and resampling units are WSIs. Patient-level quantities remain blank pending incorporation of the corresponding records.

\subsection{Hybrid-Evidence MIL}
Patch features were projected from 1,536 to 256 dimensions using a linear layer, layer normalization, GELU activation and dropout. The model constructed five vectors: gated attention pooling \citep{ref5}, the feature mean, the sample standard deviation, a high-margin feature mean, and the sum of a low-margin mean and coordinate-group summary. Two banks of eight learnable prototypes supplied maximum cosine similarities; their difference defined the margin. High- and low-margin summaries averaged the 32 extreme instances in each direction. Regional pooling operated on 16 equal-sized chunks after coordinate sorting, without an explicit anatomical-region annotation or adjacency graph.

Let $\mathbf{e}_1,\ldots,\mathbf{e}_5\in\mathbb{R}^{256}$ denote these vectors and $\mathbf{u}\in\mathbb{R}^{8}$ the similarity/margin statistics. A learned gate $g$ and fusion/classification network $f$ produced
\begin{align}
\boldsymbol\gamma&=\operatorname{softmax}\!\left(g([\mathbf{e}_1;\cdots;\mathbf{e}_5;\mathbf{u}])\right),\label{eq:main-gate}\\
p&=\operatorname{sigmoid}\!\left(f([\gamma_1\mathbf{e}_1;\cdots;\gamma_5\mathbf{e}_5;\mathbf{u}])\right).\label{eq:main-score}
\end{align}
The statistics bypassed branch scaling. The full mathematical definitions, dimensions and operators are retained in Supplementary Methods. The score $p$ is an uncalibrated MALT score. Training minimized weighted binary cross-entropy,
\begin{equation}
\mathcal L=-\frac{1}{B}\sum_{b=1}^{B}\left[\frac{77}{67}y_b\log p_b+(1-y_b)\log(1-p_b)\right],
\end{equation}
implemented from logits. The class weight was the ratio of SS to MALT training slides. No region-level or prototype-specific objective was used in the original model.

\subsection{Training and comparator analyses}
The retained Hybrid model used AdamW, learning rate $5\times10^{-5}$, weight decay $10^{-4}$, batch size eight, dropout 0.3 and 30 epochs. Seed 42 was used in the original experiment, and checkpoint selection by validation AUROC retained epoch 15. Validation performance was not interpreted as an independent test estimate. The original training environment has not been fully recovered (completion item T02).

Historical comparisons included mean pooling, ABMIL and local attention, transformer and prototype aggregators. Their architecture dimensions and training budgets differed. We use descriptive implementation names and preserve mappings to the archived code in Supplementary Table~S1. These local variants are not complete reproductions of CLAM, DSMIL, DTFD-MIL, TransMIL or R$^2$T-MIL \citep{ref7,ref14,ref15,ref16,rrt2024}. CONCH v1.5 with TITAN used a different encoder and 1,024 instances per slide and was treated as a separate pipeline \citep{ref13}. The historical three-model ensemble and its development-selected threshold are reported only as exploratory supplementary analyses.

The matched comparison trained attention-only, full fixed-scaling and full learned-gate models with seeds 17, 29 and 43. Within each seed, shared parameters, sampled training bags and minibatch orders were matched. All runs used the Hybrid optimization settings above and the first strictly maximal validation-AUROC epoch. The models differed in parameter count; the comparison did not isolate capacity from gating.

The later algorithm screen used ten configurations and the same three seeds, followed by two additional seeds (59 and 71) for the internally selected candidate and Hybrid. The six control reruns reproduced the corresponding earlier matched-attention and Hybrid runs; they were not additional independent seeds. Candidate selection used three-seed mean validation AUROC, with the complete ranking frozen before the new external inference. The selected configuration remained unchanged after the additional seeds. Normalization, attention masking, feature-subspace aggregation and two-bag consistency settings are specified in Supplementary Methods and the archived protocol \citep{acmil2024,nnmil2026,rdrop2021}.

Two linear controls fitted StandardScaler and L2-regularized logistic regression only to the 144 training slides, using each slide's full-feature mean, with or without per-patch L2 normalization. Logistic settings were $C=1$, balanced class weights, LBFGS, 3,000 maximum iterations and tolerance $10^{-8}$, without tuning. Inference on 512-patch reference means and full-feature means was reported separately.

\subsection{External evaluation, input sampling and regional outputs}
The original external predictions used the retained Hybrid checkpoint and 512 patches sampled with seed 42 plus the original manifest index. Later ABMIL and two local comparator checkpoints were evaluated on exactly those reference inputs. Source and stain archive tags were retained without converting them to confirmed histological metadata. Omitting IHC/special-stain tags was a post hoc descriptive sensitivity analysis, not a revised primary cohort. Later algorithm evaluations reused the same historical external collection after it had already been observed during development.

For input-sampling analysis, the retained Hybrid checkpoint was evaluated on 20 additional deterministic bags per slide. We calculated per-slide score standard deviations, occurrence of both classifications around 0.5, and complete-cohort AUROC for each draw. Separately, all five seeds of Hybrid and Attention+L2 were evaluated on five new bags per slide. Seed and bag ensembles used arithmetic means of probabilities. Exact sampling rules and the distinct feature-window logit averaging used by the subspace model appear in Supplementary Methods. Repeated bags and seeds were not treated as additional clinical observations.

Regional displays reused archived margin maps and patch thumbnails without altering tissue pixels. Twelve visualization cases had their performance and interpretation scores linked, because the two scripts could sample different bags. Eight displayed image windows were matched to the archived crops. These checks support image traceability; pathological regional annotation and reader comparisons remain outside the current results.

\subsection{Statistical analysis}
The primary metric was AUROC. AP was calculated by the average-precision definition; threshold metrics used 0.5 unless a supplementary analysis explicitly stated otherwise. Sensitivity, specificity, predictive values, accuracy and confusion counts were calculated from aligned predictions. Balanced accuracy was the mean of sensitivity and specificity. Brier scores were mean squared differences between scores and binary labels.

Primary fixed-checkpoint AUROC/AP intervals used 2,000 class-stratified WSI bootstrap samples with seed 20261016; proportion intervals used Wilson's method \citep{ref19}. Historical internal paired AUROC intervals used aligned resampling; later external and algorithm comparisons used 10,000 class-stratified paired samples with seed 20261104. These conditional intervals do not incorporate checkpoint/configuration selection, patient clustering or training variation. Comparisons were descriptive and not multiplicity-adjusted. Across training seeds we report individual values and means with sample standard deviations, separately from AUROC calculated after probability averaging. Statistics and figures used Python; numerical source tables and executable rendering code accompany the draft.
